\documentclass[11pt,a4paper]{article}
\usepackage[utf8]{inputenc}
\usepackage[spanish,es-minimal]{babel}
\usepackage[margin=2.5cm]{geometry}
\usepackage{amsmath,amssymb}
\usepackage{graphicx}
\usepackage{booktabs}
\usepackage{hyperref}
\usepackage{listings}
\usepackage{xcolor}
\usepackage{fancyhdr}
\usepackage{titlesec}

% Configuración de colores
\definecolor{unsblue}{RGB}{30,60,120}
\definecolor{codegray}{RGB}{240,240,240}

% Configuración de títulos
\titleformat{\section}{\Large\bfseries\color{unsblue}}{\thesection}{1em}{}
\titleformat{\subsection}{\large\bfseries\color{unsblue}}{\thesubsection}{1em}{}

% Configuración de código
\lstset{
    backgroundcolor=\color{codegray},
    basicstyle=\ttfamily\small,
    breaklines=true,
    frame=single,
    xleftmargin=1cm,
    xrightmargin=1cm,
    numbers=left,
    numberstyle=\tiny\color{gray},
    keywordstyle=\color{blue},
    commentstyle=\color{gray},
    stringstyle=\color{red}
}

% Encabezado y pie de página
\pagestyle{fancy}
\fancyhf{}
\fancyhead[L]{Métodos de Análisis de Datos 1}
\fancyhead[R]{Unidad 5}
\fancyfoot[C]{\thepage}

\title{\textbf{Métodos de Análisis de Datos 1}\\
\Large Unidad 5: Predicción y Comunicación de Resultados\\
\large Teoría Completa}
\author{Departamento de Matemática\\Universidad Nacional del Sur}
\date{Primer Semestre 2026}

\begin{document}

\maketitle

\tableofcontents
\newpage

\section{Introducción}

La Unidad 5 cierra el ciclo completo de análisis de datos abordando aspectos cruciales que conectan el trabajo técnico con su aplicación práctica y comunicación efectiva. Hasta ahora hemos aprendido a obtener, limpiar, explorar y modelar datos. En esta unidad nos enfocamos en:

\begin{itemize}
\item Llevar modelos desde el desarrollo hasta producción
\item Interpretar y explicar las decisiones de modelos complejos
\item Comunicar resultados de manera efectiva a diferentes audiencias
\item Crear reportes reproducibles y documentación de calidad
\item Considerar las implicaciones éticas del análisis de datos
\end{itemize}

Estos aspectos son frecuentemente subestimados en cursos introductorios, pero son fundamentales para el éxito profesional en ciencia de datos.

\section{Ciclo de Vida de Modelos en Producción}

\subsection{Fases del Ciclo de Vida}

El desarrollo de un modelo de machine learning no termina cuando obtenemos buenas métricas en el conjunto de validación. El ciclo completo incluye:

\subsubsection{1. Desarrollo e Investigación}

\begin{itemize}
\item \textbf{Exploración de datos:} análisis exploratorio, identificación de patrones, detección de problemas de calidad
\item \textbf{Ingeniería de features:} creación, selección y transformación de variables
\item \textbf{Experimentación:} probar diferentes algoritmos, hiperparámetros y estrategias
\item \textbf{Validación:} evaluación rigurosa con datos no vistos, análisis de errores
\end{itemize}

\textbf{Herramientas típicas:} Jupyter notebooks, VS Code, bibliotecas de ML (scikit-learn, statsmodels)

\subsubsection{2. Deployment (Puesta en Producción)}

Una vez que tenemos un modelo que cumple los requisitos de rendimiento y negocio, debemos ponerlo a disposición de usuarios o sistemas. Esto implica:

\begin{itemize}
\item \textbf{Serialización:} guardar el modelo entrenado en un formato que pueda ser cargado posteriormente
\item \textbf{Creación de API:} exponer el modelo mediante endpoints REST para que otros sistemas lo consuman
\item \textbf{Containerización:} empaquetar modelo y dependencias (Docker) para garantizar reproducibilidad
\item \textbf{Integración:} conectar el modelo con bases de datos, pipelines de datos, sistemas empresariales
\end{itemize}

\textbf{Ejemplo de flujo:}
\begin{enumerate}
\item Entrenar modelo en notebook
\item Serializar modelo con pickle/joblib
\item Crear aplicación Flask/FastAPI que carga el modelo
\item Dockerizar la aplicación
\item Deployar en servidor/cloud (AWS, Azure, GCP)
\end{enumerate}

\subsubsection{3. Monitoreo}

Los modelos en producción deben ser monitoreados continuamente para detectar:

\begin{itemize}
\item \textbf{Degradación del rendimiento:} las métricas (accuracy, AUC, etc.) pueden empeorar con el tiempo
\item \textbf{Data drift:} cambios en la distribución de las variables de entrada
\item \textbf{Concept drift:} cambios en la relación entre predictores y variable objetivo
\item \textbf{Problemas técnicos:} tiempos de respuesta, errores, disponibilidad
\end{itemize}

\textbf{Métricas a monitorear:}
\begin{itemize}
\item Performance del modelo (accuracy, precision, recall, AUC)
\item Distribución de predicciones
\item Distribución de features de entrada
\item Latencia (tiempo de respuesta)
\item Tasa de errores
\end{itemize}

\subsubsection{4. Reentrenamiento y Versionado}

Cuando detectamos degradación o disponemos de nuevos datos, debemos reentrenar:

\begin{itemize}
\item \textbf{Reentrenamiento programado:} automático cada cierto tiempo (diario, semanal, mensual)
\item \textbf{Reentrenamiento disparado:} cuando las métricas caen por debajo de un umbral
\item \textbf{Versionado de modelos:} mantener registro de versiones, poder hacer rollback
\item \textbf{A/B testing:} comparar modelo nuevo vs antiguo en producción antes de reemplazo completo
\end{itemize}

\textbf{Herramientas de versionado:} MLflow, DVC, Neptune.ai, Weights \& Biases

\subsection{Data Drift y Concept Drift}

\subsubsection{Data Drift}

Cambio en la distribución de las variables de entrada $P(X)$. Por ejemplo:

\begin{itemize}
\item Un modelo de scoring crediticio entrenado pre-pandemia puede recibir solicitudes con perfiles muy diferentes post-pandemia
\item Variables demográficas cambian con el tiempo (edad, ingresos)
\item Cambios estacionales en datos de ventas
\end{itemize}

\textbf{Detección:}
\begin{itemize}
\item Comparar distribuciones entre train y producción (test de Kolmogorov-Smirnov, chi-cuadrado)
\item Monitorear estadísticos descriptivos (media, varianza) de cada feature
\item Visualizaciones: histogramas comparativos, Q-Q plots
\end{itemize}

\subsubsection{Concept Drift}

Cambio en la relación entre predictores y objetivo $P(Y|X)$. El concepto que el modelo aprendió ya no es válido. Ejemplos:

\begin{itemize}
\item Un modelo de predicción de ventas entrenado en un mercado estable puede fallar si cambian las regulaciones
\item Detección de fraude: los estafadores adaptan sus métodos
\item Recomendaciones: las preferencias de usuarios cambian con el tiempo
\end{itemize}

\textbf{Detección:}
\begin{itemize}
\item Monitorear performance en producción (accuracy, precision, recall)
\item Comparar predicciones vs valores reales cuando estén disponibles
\item Alertas cuando métricas caen por debajo de umbrales
\end{itemize}

\subsection{Consideraciones de Deployment}

\subsubsection{Escalabilidad}

\begin{itemize}
\item \textbf{Latencia:} tiempo de respuesta debe ser aceptable (usualmente $<$ 100ms para aplicaciones interactivas)
\item \textbf{Throughput:} cantidad de predicciones por segundo
\item \textbf{Batch vs real-time:} ¿predicciones individuales o en lote?
\item \textbf{Trade-offs:} modelos más complejos pueden ser más precisos pero más lentos
\end{itemize}

\subsubsection{Mantenibilidad}

\begin{itemize}
\item \textbf{Documentación:} del modelo, features, dependencias, proceso de entrenamiento
\item \textbf{Testing:} tests unitarios para preprocesamiento, validaciones de entrada/salida
\item \textbf{Logging:} registrar predicciones, inputs, errores para debugging
\item \textbf{Reproducibilidad:} control de versiones de código, datos, modelo
\end{itemize}

\section{Serialización de Modelos}

La serialización es el proceso de convertir un objeto Python (como un modelo entrenado) en un formato que puede ser almacenado en disco y posteriormente recargado.

\subsection{Formatos de Serialización}

\subsubsection{pickle}

Formato estándar de Python para serialización de objetos.

\textbf{Ventajas:}
\begin{itemize}
\item Nativo de Python, no requiere librerías adicionales
\item Puede serializar casi cualquier objeto Python
\item Simple de usar
\end{itemize}

\textbf{Desventajas:}
\begin{itemize}
\item Específico de Python (no portable a otros lenguajes)
\item Vulnerabilidades de seguridad: no cargar pickles de fuentes no confiables
\item Puede haber incompatibilidades entre versiones de librerías
\end{itemize}

\textbf{Ejemplo:}
\begin{lstlisting}[language=Python]
import pickle
from sklearn.linear_model import LogisticRegression

# Entrenar modelo
modelo = LogisticRegression()
modelo.fit(X_train, y_train)

# Guardar
with open('modelo.pkl', 'wb') as f:
    pickle.dump(modelo, f)

# Cargar
with open('modelo.pkl', 'rb') as f:
    modelo_cargado = pickle.load(f)

# Usar
predicciones = modelo_cargado.predict(X_test)
\end{lstlisting}

\subsubsection{joblib}

Optimización de pickle para objetos grandes con arrays NumPy (común en modelos de ML).

\textbf{Ventajas:}
\begin{itemize}
\item Más eficiente que pickle para modelos grandes
\item Compresión automática
\item API simple compatible con pickle
\end{itemize}

\textbf{Ejemplo:}
\begin{lstlisting}[language=Python]
from joblib import dump, load

# Guardar
dump(modelo, 'modelo.joblib')

# Cargar
modelo_cargado = load('modelo.joblib')
\end{lstlisting}

\subsubsection{ONNX (Open Neural Network Exchange)}

Formato abierto e interoperable para modelos de ML, especialmente redes neuronales.

\textbf{Ventajas:}
\begin{itemize}
\item Portable entre frameworks (PyTorch, TensorFlow, scikit-learn)
\item Portable entre lenguajes (Python, C++, Java, JavaScript)
\item Optimizado para inferencia rápida
\end{itemize}

\textbf{Desventajas:}
\begin{itemize}
\item Más complejo de configurar
\item No todos los modelos de scikit-learn son convertibles fácilmente
\end{itemize}

\textbf{Ejemplo con scikit-learn:}
\begin{lstlisting}[language=Python]
from skl2onnx import convert_sklearn
from skl2onnx.common.data_types import FloatTensorType

# Definir tipo de input
initial_type = [('float_input', FloatTensorType([None, X_train.shape[1]]))]

# Convertir
onx = convert_sklearn(modelo, initial_types=initial_type)

# Guardar
with open("modelo.onnx", "wb") as f:
    f.write(onx.SerializeToString())
\end{lstlisting}

\subsection{Buenas Prácticas de Serialización}

\begin{enumerate}
\item \textbf{Guardar todo el pipeline:} no solo el modelo, sino también preprocesadores (scalers, encoders)
\item \textbf{Versionado:} incluir versión del modelo y de librerías en el nombre del archivo
\item \textbf{Metadata:} guardar información adicional (fecha de entrenamiento, métricas, features utilizadas)
\item \textbf{Testing post-carga:} verificar que el modelo cargado produce las mismas predicciones
\item \textbf{Compresión:} para modelos grandes, usar compresión (joblib lo hace automáticamente)
\end{enumerate}

\textbf{Ejemplo de serialización completa:}
\begin{lstlisting}[language=Python]
from sklearn.pipeline import Pipeline
from sklearn.preprocessing import StandardScaler
from sklearn.ensemble import RandomForestClassifier
from joblib import dump
import json
from datetime import datetime

# Crear pipeline completo
pipeline = Pipeline([
    ('scaler', StandardScaler()),
    ('classifier', RandomForestClassifier(n_estimators=100, random_state=42))
])

# Entrenar
pipeline.fit(X_train, y_train)

# Evaluar
accuracy = pipeline.score(X_test, y_test)

# Guardar pipeline
dump(pipeline, 'pipeline_v1.0.0.joblib')

# Guardar metadata
metadata = {
    'version': '1.0.0',
    'timestamp': datetime.now().isoformat(),
    'accuracy': accuracy,
    'features': list(X_train.columns),
    'sklearn_version': sklearn.__version__
}

with open('pipeline_v1.0.0_metadata.json', 'w') as f:
    json.dump(metadata, f, indent=2)
\end{lstlisting}

\section{Interpretabilidad de Modelos}

La interpretabilidad se refiere a la capacidad de entender y explicar cómo un modelo toma sus decisiones. Es crucial por:

\begin{itemize}
\item \textbf{Confianza:} usuarios confían más en modelos que pueden explicar
\item \textbf{Debugging:} identificar errores y sesgos en el modelo
\item \textbf{Cumplimiento regulatorio:} en algunos dominios (finanzas, salud) es obligatorio poder explicar decisiones
\item \textbf{Mejora del modelo:} entender qué aprende el modelo ayuda a mejorar features
\end{itemize}

\subsection{Modelos Interpretables vs Caja Negra}

\subsubsection{Modelos Inherentemente Interpretables}

\textbf{Regresión Lineal/Logística:}
\begin{itemize}
\item Coeficientes indican dirección y magnitud del efecto de cada variable
\item Fácil de entender: ``un aumento de 1 unidad en $X_1$ resulta en un cambio de $\beta_1$ en $Y$''
\end{itemize}

\textbf{Árboles de Decisión:}
\begin{itemize}
\item Estructura visual clara: reglas if-then
\item Fácil de explicar a no técnicos
\item Limitación: árboles profundos pierden interpretabilidad
\end{itemize}

\subsubsection{Modelos de Caja Negra}

\textbf{Random Forest, Gradient Boosting:}
\begin{itemize}
\item Combinan múltiples árboles (decenas a miles)
\item No hay una regla simple de interpretación
\item Necesitan métodos post-hoc para interpretación
\end{itemize}

\textbf{Redes Neuronales:}
\begin{itemize}
\item Millones de parámetros
\item Transformaciones no lineales complejas
\item Extremadamente difíciles de interpretar directamente
\end{itemize}

\textbf{SVM con kernels no lineales:}
\begin{itemize}
\item Operan en espacios de alta dimensión
\item Vectores de soporte no son directamente interpretables
\end{itemize}

\subsection{Trade-off Interpretabilidad vs Performance}

En general existe un trade-off:
\begin{itemize}
\item Modelos simples (regresión lineal, árboles pequeños) son fáciles de interpretar pero pueden tener menor performance
\item Modelos complejos (random forest, redes neuronales) tienen mejor performance pero son más difíciles de interpretar
\end{itemize}

\textbf{Estrategia recomendada:}
\begin{enumerate}
\item Comenzar con modelos simples e interpretables (baseline)
\item Si el performance no es suficiente, probar modelos más complejos
\item Usar técnicas de interpretabilidad (SHAP, LIME) para explicar modelos complejos
\item Documentar claramente los trade-offs
\end{enumerate}

\subsection{SHAP (SHapley Additive exPlanations)}

SHAP es un método unificado para explicar las salidas de modelos de machine learning, basado en valores de Shapley de la teoría de juegos cooperativos.

\subsubsection{Idea Principal}

Para cada predicción, SHAP calcula la contribución de cada feature al valor predicho. Formalmente:

$$f(x) = \phi_0 + \sum_{i=1}^{p} \phi_i$$

donde:
\begin{itemize}
\item $f(x)$ es la predicción del modelo para la instancia $x$
\item $\phi_0$ es el valor base (predicción promedio del modelo)
\item $\phi_i$ es la contribución de la feature $i$ (valor SHAP)
\end{itemize}

\textbf{Propiedades deseables que cumple SHAP:}
\begin{itemize}
\item \textbf{Local accuracy:} la suma de contribuciones reproduce la predicción
\item \textbf{Missingness:} features ausentes no contribuyen
\item \textbf{Consistency:} si un modelo cambia para depender más de una feature, su contribución aumenta
\end{itemize}

\subsubsection{Interpretación de SHAP Values}

\begin{itemize}
\item Valores positivos: la feature empuja la predicción hacia arriba
\item Valores negativos: la feature empuja la predicción hacia abajo
\item Magnitud: importancia de la contribución
\end{itemize}

\subsubsection{Tipos de Gráficos SHAP}

\textbf{1. Force Plot:} visualiza cómo cada feature contribuye a una predicción individual

\textbf{2. Summary Plot:} muestra la distribución de SHAP values para cada feature en todo el dataset
\begin{itemize}
\item Eje vertical: features ordenadas por importancia
\item Eje horizontal: valor SHAP (impacto en predicción)
\item Color: valor de la feature (alto/bajo)
\end{itemize}

\textbf{3. Dependence Plot:} muestra la relación entre el valor de una feature y su SHAP value

\textbf{4. Waterfall Plot:} descomposición aditiva de una predicción individual

\subsubsection{Uso en Python}

\begin{lstlisting}[language=Python]
import shap
from sklearn.ensemble import RandomForestClassifier

# Entrenar modelo
modelo = RandomForestClassifier()
modelo.fit(X_train, y_train)

# Crear explainer
explainer = shap.TreeExplainer(modelo)

# Calcular SHAP values
shap_values = explainer.shap_values(X_test)

# Visualizar
# Summary plot (global)
shap.summary_plot(shap_values[1], X_test)

# Force plot (individual, primera instancia)
shap.force_plot(explainer.expected_value[1], 
                shap_values[1][0], 
                X_test.iloc[0])

# Dependence plot
shap.dependence_plot("feature_name", shap_values[1], X_test)
\end{lstlisting}

\subsection{LIME (Local Interpretable Model-agnostic Explanations)}

LIME explica predicciones individuales aproximando localmente el modelo con un modelo interpretable simple.

\subsubsection{Idea Principal}

Para explicar una predicción:
\begin{enumerate}
\item Generar perturbaciones (variaciones) del punto a explicar
\item Obtener predicciones del modelo original para estas perturbaciones
\item Entrenar un modelo interpretable simple (e.g., regresión lineal) en las perturbaciones, ponderando más las cercanas al punto original
\item El modelo simple aproxima localmente el comportamiento del modelo complejo
\end{enumerate}

\textbf{Analogía:} es como aproximar una función compleja no lineal con una recta tangente en un punto específico.

\subsubsection{Ventajas de LIME}

\begin{itemize}
\item \textbf{Model-agnostic:} funciona con cualquier modelo de caja negra
\item \textbf{Explicaciones locales:} se enfoca en una predicción específica
\item \textbf{Interpretable:} usa modelos lineales como explicadores
\end{itemize}

\subsubsection{Desventajas de LIME}

\begin{itemize}
\item \textbf{Inestabilidad:} pequeños cambios en el punto pueden cambiar mucho la explicación
\item \textbf{Dependencia de hiperparámetros:} tamaño de vecindad, kernel
\item \textbf{No garantiza propiedades globales:} solo válido localmente
\end{itemize}

\subsubsection{Uso en Python}

\begin{lstlisting}[language=Python]
import lime
import lime.lime_tabular

# Crear explainer
explainer = lime.lime_tabular.LimeTabularExplainer(
    X_train.values,
    feature_names=X_train.columns,
    class_names=['Clase 0', 'Clase 1'],
    mode='classification'
)

# Explicar una instancia
i = 0  # índice de la instancia
exp = explainer.explain_instance(
    X_test.values[i], 
    modelo.predict_proba,
    num_features=10
)

# Visualizar
exp.show_in_notebook()

# O como gráfico
exp.as_pyplot_figure()
\end{lstlisting}

\subsection{Partial Dependence Plots (PDP)}

Los PDP muestran el efecto marginal de una feature sobre la predicción, manteniendo las demás constantes.

\subsubsection{Cálculo}

Para una feature $X_j$:
$$PD_{X_j}(x_j) = \mathbb{E}_{X_{-j}}[f(x_j, X_{-j})] = \frac{1}{n}\sum_{i=1}^{n}f(x_j, x_{-j}^{(i)})$$

Es decir, fijamos $X_j$ en un valor $x_j$, promediamos las predicciones sobre todos los valores observados de las demás features.

\subsubsection{Interpretación}

\begin{itemize}
\item Eje horizontal: valores de la feature
\item Eje vertical: predicción promedio
\item La curva muestra cómo cambia la predicción al variar la feature
\item Pendiente positiva: relación positiva
\item Pendiente negativa: relación negativa
\item Curva plana: la feature no afecta la predicción
\end{itemize}

\subsubsection{Limitaciones}

\begin{itemize}
\item \textbf{Supuesto de independencia:} asume que las features son independientes (raramente cierto)
\item \textbf{No captura interacciones:} promedia sobre todas las combinaciones
\item \textbf{Puede ser engañoso:} si hay interacciones fuertes
\end{itemize}

\textbf{Solución para interacciones:} Individual Conditional Expectation (ICE) plots, que muestran el efecto para cada instancia individualmente.

\subsubsection{Uso en Python}

\begin{lstlisting}[language=Python]
from sklearn.inspection import PartialDependenceDisplay

# PDP para una feature
fig, ax = plt.subplots(figsize=(10, 6))
PartialDependenceDisplay.from_estimator(
    modelo, X_train, 
    features=['feature_name'],
    ax=ax
)
plt.show()

# PDP para múltiples features
fig, ax = plt.subplots(figsize=(12, 4))
PartialDependenceDisplay.from_estimator(
    modelo, X_train, 
    features=['feature1', 'feature2', ('feature1', 'feature2')],
    ax=ax
)
plt.show()
\end{lstlisting}

\subsection{Feature Importance}

Mide la importancia relativa de cada feature para el modelo.

\subsubsection{Importancia en Árboles y Random Forest}

Para modelos basados en árboles, la importancia se calcula como la disminución promedio en impureza (Gini o entropía) causada por cada feature.

\textbf{Ventajas:}
\begin{itemize}
\item Rápida de calcular
\item Incorporada en scikit-learn
\end{itemize}

\textbf{Desventajas:}
\begin{itemize}
\item Sesgada hacia features con alta cardinalidad
\item No distingue correlación de causalidad
\end{itemize}

\begin{lstlisting}[language=Python]
# Obtener importancias
importances = modelo.feature_importances_

# Visualizar
indices = np.argsort(importances)[::-1]
plt.figure(figsize=(10, 6))
plt.bar(range(X_train.shape[1]), importances[indices])
plt.xticks(range(X_train.shape[1]), 
           X_train.columns[indices], 
           rotation=90)
plt.title("Feature Importances")
plt.tight_layout()
plt.show()
\end{lstlisting}

\subsubsection{Permutation Importance}

Mide el aumento en el error del modelo cuando se permutan aleatoriamente los valores de una feature (rompiendo su relación con el objetivo).

\textbf{Ventajas:}
\begin{itemize}
\item Model-agnostic
\item No sesgada por cardinalidad
\item Refleja importancia real para predicción
\end{itemize}

\textbf{Desventajas:}
\begin{itemize}
\item Más lenta de calcular
\item Puede ser confusa con features correlacionadas
\end{itemize}

\begin{lstlisting}[language=Python]
from sklearn.inspection import permutation_importance

# Calcular
perm_importance = permutation_importance(
    modelo, X_test, y_test,
    n_repeats=10, random_state=42
)

# Visualizar
sorted_idx = perm_importance.importances_mean.argsort()[::-1]
plt.figure(figsize=(10, 6))
plt.boxplot(perm_importance.importances[sorted_idx].T,
            labels=X_test.columns[sorted_idx],
            vert=False)
plt.xlabel("Decrease in accuracy")
plt.tight_layout()
plt.show()
\end{lstlisting}

\section{Comunicación de Resultados}

La comunicación efectiva es tan importante como el análisis técnico. Los mejores modelos son inútiles si no podemos comunicar sus hallazgos de manera clara y convincente.

\subsection{Principios de Comunicación Efectiva}

\subsubsection{1. Conocer a la Audiencia}

\textbf{Audiencia técnica (científicos de datos, ingenieros):}
\begin{itemize}
\item Puede incluir detalles técnicos (hiperparámetros, métricas específicas)
\item Aprecian gráficos de diagnóstico (residuos, curvas ROC)
\item Interesados en reproducibilidad y metodología
\end{itemize}

\textbf{Audiencia de negocio (gerentes, ejecutivos):}
\begin{itemize}
\item Enfoque en impacto, ROI, decisiones accionables
\item Visualizaciones simples y claras
\item Lenguaje no técnico, evitar jerga
\item Destacar resultados clave y recomendaciones
\end{itemize}

\textbf{Audiencia mixta:}
\begin{itemize}
\item Resumen ejecutivo al inicio
\item Detalles técnicos en apéndices
\item Gráficos intuitivos con explicaciones breves
\end{itemize}

\subsubsection{2. Estructura Clara}

\textbf{Reporte típico:}
\begin{enumerate}
\item \textbf{Resumen ejecutivo (1 página):} problema, metodología, resultados clave, recomendaciones
\item \textbf{Introducción:} contexto, objetivos, preguntas de investigación
\item \textbf{Datos:} fuentes, tamaño, calidad, limitaciones
\item \textbf{Metodología:} análisis exploratorio, features, modelos probados, validación
\item \textbf{Resultados:} métricas, comparaciones, hallazgos principales
\item \textbf{Conclusiones:} interpretación, implicaciones, recomendaciones
\item \textbf{Limitaciones:} supuestos, sesgos, incertidumbre
\item \textbf{Próximos pasos:} mejoras futuras, áreas de investigación
\item \textbf{Apéndices:} detalles técnicos, código, tablas completas
\end{enumerate}

\subsubsection{3. Visualizaciones Efectivas}

\textbf{Principios:}
\begin{itemize}
\item \textbf{Simplicidad:} un gráfico, un mensaje
\item \textbf{Claridad:} ejes etiquetados, leyendas, títulos descriptivos
\item \textbf{Honestidad:} no distorsionar (escalas apropiadas, barras desde cero)
\item \textbf{Accesibilidad:} colores apropiados para daltónicos, tamaño de fuente legible
\item \textbf{Consistencia:} mismo estilo a lo largo del reporte
\end{itemize}

\textbf{Elegir el gráfico apropiado:}
\begin{itemize}
\item Comparación de categorías: barras, columnas
\item Evolución temporal: líneas
\item Distribución: histogramas, boxplots, violines
\item Relación entre dos variables: scatter plots
\item Proporciones: pie charts (usar con moderación), treemaps
\item Múltiples variables: heatmaps, parallel coordinates
\end{itemize}

\subsubsection{4. Narrar una Historia}

Los datos por sí solos no son convincentes. Necesitamos construir una narrativa:

\begin{itemize}
\item \textbf{Contexto:} ¿por qué es importante este análisis?
\item \textbf{Problema:} ¿qué pregunta estamos respondiendo?
\item \textbf{Viaje:} ¿qué descubrimos en el camino?
\item \textbf{Clímax:} hallazgo principal
\item \textbf{Resolución:} ¿qué debemos hacer con esta información?
\end{itemize}

\textbf{Técnicas narrativas:}
\begin{itemize}
\item Comenzar con un ejemplo concreto o anécdota
\item Usar analogías para conceptos complejos
\item Anticipar y responder objeciones
\item Terminar con un llamado a la acción claro
\end{itemize}

\subsection{Reportar Incertidumbre}

Todo modelo tiene incertidumbre. Es crucial comunicarla honestamente:

\begin{itemize}
\item \textbf{Intervalos de confianza:} no solo puntos estimados
\item \textbf{Errores de predicción:} RMSE, MAE contextualizados
\item \textbf{Limitaciones de los datos:} sesgos de muestreo, datos faltantes
\item \textbf{Supuestos del modelo:} qué asumimos y qué tan razonables son
\end{itemize}

\textbf{Ejemplo de buena práctica:}

``El modelo predice una tasa de churn del 15\% con un intervalo de confianza del 95\% de [12\%, 18\%]. Esta estimación asume que el comportamiento futuro de los clientes será similar al histórico. Cambios en el mercado o en nuestros productos podrían afectar esta predicción.''

\subsection{Lenguaje Claro y Preciso}

\textbf{Evitar:}
\begin{itemize}
\item Jerga innecesaria (``hiperplano de máximo margen'')
\item Siglas no definidas (AUC, RMSE sin explicar)
\item Lenguaje vago (``el modelo funciona bien'')
\item Causalidad cuando solo hay correlación (``X causa Y'')
\end{itemize}

\textbf{Preferir:}
\begin{itemize}
\item Términos simples y descriptivos
\item Definiciones breves en primera mención
\item Cuantificación precisa (``el modelo clasifica correctamente el 85\% de los casos'')
\item Lenguaje asociativo (``X está asociado con Y'', ``X predice Y'')
\end{itemize}

\section{Reportes Reproducibles}

La reproducibilidad es fundamental en ciencia de datos. Un análisis reproducible puede ser ejecutado por otra persona y obtener los mismos resultados.

\subsection{Jupyter Notebooks}

Jupyter combina código, visualizaciones y narrativa en un solo documento interactivo.

\subsubsection{Ventajas}

\begin{itemize}
\item \textbf{Interactividad:} ejecutar código celda por celda
\item \textbf{Integración:} código, texto (Markdown), ecuaciones (LaTeX), imágenes
\item \textbf{Compartible:} formato JSON, puede versionarse con git
\item \textbf{Exportable:} HTML, PDF, presentaciones
\item \textbf{Popular:} ampliamente usado en industria y academia
\end{itemize}

\subsubsection{Buenas Prácticas}

\begin{enumerate}
\item \textbf{Orden lógico:} celdas ejecutables de arriba a abajo secuencialmente
\item \textbf{Restart \& Run All:} asegurarse que el notebook funcione desde cero
\item \textbf{Markdown generoso:} explicar qué hace cada sección
\item \textbf{Visualizaciones inline:} configurar \texttt{\%matplotlib inline}
\item \textbf{Control de versiones:} hacer commits frecuentes
\item \textbf{Evitar outputs grandes:} limpiar antes de commitear
\item \textbf{Gestión de dependencias:} incluir requirements.txt o environment.yml
\end{enumerate}

\subsubsection{Estructura Recomendada}

\begin{enumerate}
\item \textbf{Título y resumen:} qué hace este notebook
\item \textbf{Imports:} todas las librerías al inicio
\item \textbf{Configuración:} seeds, parámetros, paths
\item \textbf{Carga de datos:} con verificaciones básicas
\item \textbf{Exploración:} análisis exploratorio
\item \textbf{Preprocesamiento:} limpieza, transformaciones
\item \textbf{Modelado:} entrenamiento, validación
\item \textbf{Evaluación:} métricas, gráficos
\item \textbf{Conclusiones:} hallazgos principales
\item \textbf{Próximos pasos:} mejoras futuras
\end{enumerate}

\subsubsection{Exportación}

\textbf{HTML (recomendado para compartir):}
\begin{lstlisting}[language=bash]
jupyter nbconvert --to html notebook.ipynb
\end{lstlisting}

\textbf{PDF (requiere LaTeX instalado):}
\begin{lstlisting}[language=bash]
jupyter nbconvert --to pdf notebook.ipynb
\end{lstlisting}

\textbf{Presentación (slides con reveal.js):}
\begin{lstlisting}[language=bash]
jupyter nbconvert --to slides notebook.ipynb --post serve
\end{lstlisting}

\subsection{R Markdown}

R Markdown permite combinar código R, narrativa y análisis en documentos reproducibles.

\subsubsection{Estructura}

Un archivo .Rmd tiene tres componentes:

\textbf{1. Encabezado YAML:} metadatos y opciones
\begin{lstlisting}
---
title: "Análisis de Ventas Q1 2026"
author: "José Bavio"
date: "2026-03-05"
output: 
  html_document:
    toc: true
    toc_float: true
    theme: cosmo
---
\end{lstlisting}

\textbf{2. Texto Markdown:} narrativa, explicaciones

\textbf{3. Chunks de código R:} entre \texttt{\`{}\`{}\`\{r\}} y \texttt{\`{}\`{}\`}

\subsubsection{Opciones de Chunks}

Controlan cómo se ejecuta y muestra el código:

\begin{itemize}
\item \texttt{echo=FALSE}: no mostrar código, solo output
\item \texttt{eval=FALSE}: mostrar código pero no ejecutar
\item \texttt{include=FALSE}: ejecutar pero no mostrar nada
\item \texttt{message=FALSE, warning=FALSE}: suprimir mensajes
\item \texttt{fig.width=8, fig.height=6}: tamaño de gráficos
\item \texttt{cache=TRUE}: cachear resultados (útil para cálculos lentos)
\end{itemize}

\subsubsection{Formatos de Salida}

\textbf{HTML (interactivo, con gráficos dinámicos):}
\begin{lstlisting}
output: html_document
\end{lstlisting}

\textbf{PDF (vía LaTeX):}
\begin{lstlisting}
output: pdf_document
\end{lstlisting}

\textbf{Word:}
\begin{lstlisting}
output: word_document
\end{lstlisting}

\textbf{Presentaciones:}
\begin{lstlisting}
output: 
  ioslides_presentation  # Google IO style
  slidy_presentation     # W3C style
  beamer_presentation    # LaTeX Beamer
\end{lstlisting}

\subsubsection{Parámetros}

R Markdown permite crear reportes parametrizados:

\begin{lstlisting}
---
title: "Reporte de Ventas"
params:
  region: "Norte"
  año: 2026
---

## Ventas en `r params$region` - `r params$año`

```{r}
datos_filtrados <- datos %>%
  filter(region == params$region, año == params$año)
```
\end{lstlisting}

Renderizar con parámetros diferentes:
\begin{lstlisting}[language=R]
rmarkdown::render("reporte.Rmd", 
                  params = list(region = "Sur", año = 2025))
\end{lstlisting}

Esto permite automatizar reportes para diferentes regiones, períodos, etc.

\subsection{Control de Versiones}

Git es esencial para reproducibilidad:

\begin{itemize}
\item \textbf{Historial completo:} registra todos los cambios
\item \textbf{Colaboración:} múltiples personas trabajando simultáneamente
\item \textbf{Branching:} experimentar sin afectar código principal
\item \textbf{Rollback:} volver a versiones anteriores
\end{itemize}

\textbf{Workflow básico:}
\begin{lstlisting}[language=bash]
# Inicializar repositorio
git init

# Agregar archivos
git add notebook.ipynb

# Commit
git commit -m "Agregar análisis exploratorio"

# Subir a remoto (GitHub, GitLab)
git push origin main
\end{lstlisting}

\textbf{.gitignore para data science:}
\begin{lstlisting}
# Datos
data/
*.csv
*.xlsx

# Modelos entrenados
*.pkl
*.joblib
*.h5

# Notebooks con outputs
.ipynb_checkpoints/

# Entornos virtuales
venv/
env/

# IDE
.vscode/
.idea/
\end{lstlisting}

\subsection{Gestión de Dependencias}

Documentar qué librerías y versiones se usaron:

\textbf{Python - requirements.txt:}
\begin{lstlisting}
pandas==2.0.1
numpy==1.24.3
scikit-learn==1.2.2
matplotlib==3.7.1
seaborn==0.12.2
\end{lstlisting}

Generar automáticamente:
\begin{lstlisting}[language=bash]
pip freeze > requirements.txt
\end{lstlisting}

Instalar desde requirements:
\begin{lstlisting}[language=bash]
pip install -r requirements.txt
\end{lstlisting}

\textbf{Python - conda environment.yml:}
\begin{lstlisting}
name: mad1
channels:
  - conda-forge
  - defaults
dependencies:
  - python=3.10
  - pandas=2.0.1
  - scikit-learn=1.2.2
  - matplotlib=3.7.1
\end{lstlisting}

\textbf{R - renv:}
\begin{lstlisting}[language=R]
# Inicializar
renv::init()

# Snapshot (guardar estado actual)
renv::snapshot()

# Restore (reinstalar desde snapshot)
renv::restore()
\end{lstlisting}

\section{Ética en Análisis de Datos}

El análisis de datos tiene poder para afectar vidas. Con este poder viene responsabilidad.

\subsection{Sesgo en Datos y Modelos}

\subsubsection{Fuentes de Sesgo}

\textbf{1. Sesgo de Selección:}
\begin{itemize}
\item Datos no representativos de la población objetivo
\item Ejemplo: entrenar modelo de contratación con datos históricos sesgados
\end{itemize}

\textbf{2. Sesgo de Medición:}
\begin{itemize}
\item Errores sistemáticos en cómo se recolectan los datos
\item Ejemplo: encuestas telefónicas excluyen a quienes no tienen teléfono
\end{itemize}

\textbf{3. Sesgo Algorítmico:}
\begin{itemize}
\item El algoritmo amplifica sesgos existentes
\item Ejemplo: sistema de recomendación que refuerza burbujas de filtro
\end{itemize}

\textbf{4. Sesgo de Confirmación:}
\begin{itemize}
\item Analistas buscan/interpretan datos de manera que confirme sus creencias
\item Ejemplo: cherry-picking resultados favorables
\end{itemize}

\subsubsection{Consecuencias del Sesgo}

\begin{itemize}
\item \textbf{Discriminación:} decisiones injustas contra grupos protegidos
\item \textbf{Profecía autocumplida:} modelo sesgado perpetúa desigualdad
\item \textbf{Pérdida de confianza:} en sistemas algorítmicos
\item \textbf{Daño legal:} violación de leyes anti-discriminación
\end{itemize}

\textbf{Caso de estudio:} COMPAS (Correctional Offender Management Profiling for Alternative Sanctions)

Un sistema de evaluación de riesgo de reincidencia usado en el sistema judicial de EEUU fue encontrado con sesgo racial: predecía falsamente alta reincidencia para acusados afroamericanos.

\subsubsection{Mitigación de Sesgo}

\begin{enumerate}
\item \textbf{Auditoría de datos:} 
\begin{itemize}
\item ¿Los datos son representativos?
\item ¿Hay subrepresentación de algún grupo?
\item ¿Cómo se recolectaron los datos?
\end{itemize}

\item \textbf{Análisis de equidad (fairness):}
\begin{itemize}
\item Métricas de equidad: disparate impact, equalized odds
\item Evaluar performance por subgrupo (género, raza, edad)
\item Detectar diferencias sistemáticas
\end{itemize}

\item \textbf{Técnicas de mitigación:}
\begin{itemize}
\item Pre-procesamiento: rebalancear datos, eliminar variables sensibles
\item In-processing: regularización de fairness, adversarial debiasing
\item Post-procesamiento: ajustar umbrales por grupo
\end{itemize}

\item \textbf{Diversidad en equipos:}
\begin{itemize}
\item Perspectivas diversas ayudan a identificar sesgos
\item Incluir stakeholders afectados en el proceso
\end{itemize}

\item \textbf{Transparencia y documentación:}
\begin{itemize}
\item Documentar limitaciones conocidas
\item Datasheets for datasets, Model Cards
\end{itemize}
\end{enumerate}

\subsection{Fairness y Discriminación Algorítmica}

Existen múltiples definiciones de fairness, a menudo incompatibles entre sí:

\subsubsection{1. Demographic Parity (Paridad Demográfica)}

El modelo predice la misma proporción de resultados positivos para todos los grupos:
$$P(\hat{Y}=1 | A=a) = P(\hat{Y}=1 | A=b)$$

donde $A$ es el atributo sensible (raza, género).

\textbf{Ejemplo:} mismo porcentaje de hombres y mujeres aprobados para un crédito.

\subsubsection{2. Equalized Odds}

Las tasas de verdaderos positivos y falsos positivos son iguales entre grupos:
$$P(\hat{Y}=1 | Y=1, A=a) = P(\hat{Y}=1 | Y=1, A=b)$$
$$P(\hat{Y}=1 | Y=0, A=a) = P(\hat{Y}=1 | Y=0, A=b)$$

\textbf{Ejemplo:} el modelo identifica correctamente a candidatos cualificados y no cualificados por igual en ambos grupos.

\subsubsection{3. Predictive Parity}

El valor predictivo positivo es igual entre grupos:
$$P(Y=1 | \hat{Y}=1, A=a) = P(Y=1 | \hat{Y}=1, A=b)$$

\textbf{Ejemplo:} entre quienes el modelo predice reincidencia, el porcentaje que realmente reincide es el mismo en ambos grupos.

\subsubsection{Tensiones entre Definiciones}

En general, no es posible satisfacer múltiples definiciones simultáneamente. Debemos elegir basándonos en el contexto:

\begin{itemize}
\item \textbf{Contratación:} demographic parity puede ser apropiada (cuotas)
\item \textbf{Justicia penal:} equalized odds (evitar discriminación en precisión)
\item \textbf{Medicina:} predictive parity (confianza en diagnósticos)
\end{itemize}

\subsection{Privacidad y Protección de Datos}

\subsubsection{Principios}

\begin{enumerate}
\item \textbf{Minimización de datos:} recolectar solo lo necesario
\item \textbf{Propósito específico:} usar datos solo para lo autorizado
\item \textbf{Consentimiento informado:} las personas deben saber qué se recolecta y cómo se usa
\item \textbf{Transparencia:} derecho a saber qué datos se tienen
\item \textbf{Derecho al olvido:} poder eliminar datos personales
\item \textbf{Seguridad:} proteger datos contra acceso no autorizado
\end{enumerate}

\subsubsection{Regulaciones}

\textbf{GDPR (General Data Protection Regulation, Unión Europea):}
\begin{itemize}
\item Protección de datos personales de ciudadanos UE
\item Consentimiento explícito requerido
\item Derecho a explicación de decisiones algorítmicas
\item Multas severas por incumplimiento
\end{itemize}

\textbf{CCPA (California Consumer Privacy Act):}
\begin{itemize}
\item Derechos de consumidores sobre sus datos
\item Opt-out de venta de datos
\item Transparencia en recolección y uso
\end{itemize}

\textbf{Argentina - Ley 25.326 de Protección de Datos Personales:}
\begin{itemize}
\item Regula tratamiento de datos personales
\item Consentimiento del titular
\item Derecho de acceso, rectificación y supresión
\end{itemize}

\subsubsection{Técnicas de Privacidad}

\textbf{1. Anonimización:}
\begin{itemize}
\item Eliminar identificadores directos (nombre, DNI, dirección)
\item Problema: reidentificación mediante quasi-identificadores
\end{itemize}

\textbf{2. K-anonimity:}
\begin{itemize}
\item Cada registro es indistinguible de al menos $k-1$ otros
\item Generalización y supresión de atributos
\end{itemize}

\textbf{3. Differential Privacy:}
\begin{itemize}
\item Agregar ruido aleatorio a los resultados
\item Garantía matemática de privacidad individual
\item Trade-off entre privacidad y utilidad
\end{itemize}

\textbf{4. Federated Learning:}
\begin{itemize}
\item Entrenar modelos sin centralizar datos
\item Datos permanecen en dispositivos locales
\item Solo se comparten actualizaciones del modelo
\end{itemize}

\subsection{Responsabilidad en Uso de Modelos Predictivos}

\subsubsection{Limitaciones y Comunicación}

\begin{itemize}
\item \textbf{Documentar supuestos:} qué condiciones deben cumplirse para que el modelo funcione
\item \textbf{Reportar incertidumbre:} intervalos de confianza, márgenes de error
\item \textbf{Advertir sobre mal uso:} para qué NO debe usarse el modelo
\item \textbf{Actualización periódica:} modelos se vuelven obsoletos
\end{itemize}

\subsubsection{Decisiones Humanas vs Algorítmicas}

\textbf{Cuándo usar algoritmos:}
\begin{itemize}
\item Decisiones repetitivas, alto volumen
\item Criterios objetivos claramente definidos
\item Datos de alta calidad disponibles
\item Consecuencias moderadas de errores
\end{itemize}

\textbf{Cuándo mantener humanos en el loop:}
\begin{itemize}
\item Decisiones con alto impacto individual (salud, justicia)
\item Contextos complejos que requieren juicio
\item Situaciones novedosas no vistas en entrenamiento
\item Casos donde explicación y empatía son importantes
\end{itemize}

\textbf{Enfoque recomendado:} ``Human-in-the-loop''
\begin{itemize}
\item Modelo hace recomendación
\item Humano toma decisión final
\item Humano puede anular al modelo con justificación
\item Feedback de decisiones humanas mejora modelo
\end{itemize}

\subsubsection{Auditoría y Rendición de Cuentas}

\begin{itemize}
\item \textbf{Auditorías regulares:} revisar performance, sesgo, errores
\item \textbf{Trazabilidad:} registrar decisiones y sus justificaciones
\item \textbf{Responsables identificados:} quién responde por el sistema
\item \textbf{Mecanismos de apelación:} las personas deben poder disputar decisiones
\item \textbf{Evaluaciones de impacto:} antes del despliegue, evaluar riesgos
\end{itemize}

\subsection{Consideraciones Éticas en Proyectos}

\subsubsection{Preguntas a Hacerse}

Antes de empezar un proyecto de análisis de datos:

\begin{enumerate}
\item \textbf{¿Tenemos derecho a estos datos?} ¿Hubo consentimiento informado?
\item \textbf{¿Para qué se usarán los resultados?} ¿Podría usarse para propósitos dañinos?
\item \textbf{¿Quién se beneficia y quién podría ser perjudicado?}
\item \textbf{¿Los datos son representativos?} ¿Hay sesgos conocidos?
\item \textbf{¿Qué tan confiables son las predicciones?} ¿Cuál es el margen de error?
\item \textbf{¿Hay alternativas no algorítmicas?} ¿Son preferibles?
\item \textbf{¿Quién será responsable si algo sale mal?}
\item \textbf{¿Podemos explicar las decisiones?} ¿A quién, cómo?
\end{enumerate}

\subsubsection{Red Flags Éticas}

Señales de alerta que indican problemas éticos potenciales:

\begin{itemize}
\item Cliente solicita ``eliminar variables sensibles pero mantener performance''
\item Objetivos vagos o cambiantes (``solo queremos ver qué encontramos'')
\item Presión para reportar solo resultados favorables
\item Falta de transparencia sobre uso final
\item Datos obtenidos de manera cuestionable
\item Stakeholders afectados no consultados
\item Consecuencias de error asimétricas (falsos negativos mucho peores que falsos positivos, o viceversa)
\end{itemize}

\subsection{Código de Ética}

Organizaciones profesionales han desarrollado códigos de ética:

\textbf{ACM Code of Ethics (Association for Computing Machinery):}
\begin{itemize}
\item Contribuir al bienestar humano y social
\item Evitar daño
\item Ser honesto y confiable
\item Ser justo y no discriminar
\item Respetar privacidad
\item Honrar confidencialidad
\end{itemize}

\textbf{Principios a seguir como científicos de datos:}
\begin{enumerate}
\item \textbf{Integridad:} reportar hallazgos honestamente, incluso si no son favorables
\item \textbf{Transparencia:} documentar limitaciones, supuestos, incertidumbres
\item \textbf{Responsabilidad:} considerar impacto potencial de nuestro trabajo
\item \textbf{Equidad:} esforzarse por crear sistemas justos y no discriminatorios
\item \textbf{Privacidad:} proteger datos personales y respetar consentimiento
\item \textbf{Competencia:} reconocer límites de nuestra expertise, buscar ayuda cuando sea necesario
\item \textbf{Beneficencia:} buscar que nuestro trabajo beneficie a la sociedad
\end{enumerate}

\section{Resumen}

La Unidad 5 completa el ciclo de análisis de datos enfocándose en aspectos cruciales pero frecuentemente descuidados:

\begin{itemize}
\item \textbf{Ciclo de vida en producción:} llevar modelos desde desarrollo hasta deployment, monitoreo y reentrenamiento
\item \textbf{Serialización:} guardar y cargar modelos con pickle, joblib, ONNX
\item \textbf{Interpretabilidad:} explicar decisiones de modelos complejos con SHAP, LIME, PDP
\item \textbf{Comunicación:} transmitir hallazgos de manera efectiva a audiencias técnicas y no técnicas
\item \textbf{Reportes reproducibles:} Jupyter notebooks y R Markdown para documentación ejecutable
\item \textbf{Ética:} reconocer y mitigar sesgos, proteger privacidad, ser responsables en el uso de datos
\end{itemize}

Estos aspectos son fundamentales para el éxito profesional. Un modelo técnicamente perfecto es inútil si no puede desplegarse, explicarse o comunicarse adecuadamente. Y un análisis brillante puede causar daño si no se consideran sus implicaciones éticas.

Como científicos de datos, tenemos la responsabilidad de:
\begin{itemize}
\item Crear modelos que funcionen en el mundo real, no solo en notebooks
\item Explicar nuestros hallazgos de manera clara y honesta
\item Documentar nuestro trabajo para que otros puedan reproducirlo y auditarlo
\item Considerar el impacto de nuestras decisiones técnicas en las personas
\item Abogar por el uso ético de datos y algoritmos
\end{itemize}

La ciencia de datos es poder. Usemos ese poder sabiamente.

\end{document}
